% GENERATED FILE. Edit canonical lessons, then run npm run generate:books.
% canonical-source-hash: fnv1a64:9016d93cc30caf80
% canonical-lessons: ZH-C283-zenme, ZH-W283-zen, ZH-C283-zenme-zou, ZH-W283-zou, ZH-C283-wang, ZH-C283-zai-again, ZH-C283-luguo, ZH-C283-practice

\chapter{Which Way?}
\label{ch:zh-which-way}
\hlchaptermodality{283}

\begin{chapteropening}
\textbf{By the end of this chapter:} \emph{I can ask the way and follow simple directions.}

Asking the way with \zh{去} … \zh{怎么走}, following directions with wǎng, saying again with \zh{再} and passing by on the way, built on two new written characters: \zh{怎} and \zh{走}.
\end{chapteropening}

\section[\texorpdfstring{\zh{怎么} (zěnme)}{怎么 (zěnme)}]{\zh{怎么} (zěnme) — how}

\label{lesson:ZH-C283-zenme}



\begin{quote}
A new chapter, and it opens on the word you ask a stranger with.

\begin{itemize}
\raggedright
  \item \emph{Recall:} say \zh{你应该去看医生} — \emph{nǐ yīnggāi qù kàn yīshēng} — \textbf{R1}, the chapter you have finished
  \item \emph{Recall:} say \emph{zěnmeyàng}, how is it — \textbf{R3}, durable
  \item \emph{Recall:} say \emph{zěnme bàn}, what should I do — \textbf{R3}
\end{itemize}
\end{quote}

\subsection*{You'll want to know: \zh{怎么}}
\textbf{\zh{怎么}} — \emph{zěnme} — \textquotedblleft{}how; in what way\textquotedblright{}.

\begin{quote}
\textbf{\zh{这个字怎么写}} \emph{zhège zì zěnme xiě} — How do you write this character?
\end{quote}

\begin{quote}
\textbf{\zh{汉语怎么说}} \emph{Hànyǔ zěnme shuō} — How do you say it in Chinese?
\end{quote}

You have said it inside \emph{zěnmeyàng} and \emph{zěnme bàn}; here it stands alone, before a verb. The next lesson puts the first character in your hand.

\begin{cousinweb}
\textbf{\zh{怎么}} (\emph{zěnme}) is built from \zh{怎} 'how' and the ending \zh{么} of \zh{什么}. No dictionary we use records an older form of either, so its history is not settled.
\end{cousinweb}

\begin{grammarlens}[title={What you've built}]
\zh{怎么} goes right before the verb it asks about: \zh{怎么写}, \zh{怎么说}.

The answer shows the way, or gives it: \zh{这么写}, written like this.
\end{grammarlens}

\subsection*{Your turn}
\begin{itemize}
\raggedright
  \item \emph{Say it:} \emph{zěnme}
  \item \emph{Say it:} \emph{zěnme}, once more
  \item \emph{Say it:} \zh{这个字怎么写} — \emph{zhège zì zěnme xiě} — How do you write this character?
  \item \emph{Say it:} \zh{汉语怎么说} — \emph{Hànyǔ zěnme shuō} — How do you say it in Chinese?
  \item \emph{Recall:} say \emph{zěnmeyàng}, then \zh{怎么}: how is it, then how
\end{itemize}

\begin{tcolorbox}[breakable,colback=teal!4,colframe=teal!35!black,title={Before you move on}]
What does \zh{怎么} mean? (\textquotedblleft{}How\textquotedblright{}.) Where does it go? (\textbf{Right before the verb}.)
\end{tcolorbox}
\section[\texorpdfstring{\zh{怎} (zěn)}{怎 (zěn)}]{\zh{怎} — the right of \zh{作}, then \zh{心}}

\label{lesson:ZH-W283-zen}



\begin{quote}
The word first, then the character inside it.

\begin{itemize}
\raggedright
  \item \emph{Recall:} say \emph{zěnme} — \textbf{R1}, one lesson back, before it decays
\end{itemize}

\begin{quote}
\zh{作} \zh{心}
\end{quote}

\emph{Read it:} each one aloud, with nothing said first — \emph{zuò}, \emph{xīn}
\end{quote}

\subsection*{Script — the right of \zh{作}, then \zh{心}}
\hlblockfigure{\detokenize{figures/ZH-W283-zen-filmstrip.pdf}}{How \zh{怎} is written, stroke by stroke}

\begin{quote}
\zh{怎}
\end{quote}

\textbf{Two pieces, and you write both already.}

\textbf{The right of \zh{作}, five strokes}, here on top: a short left-falling stroke; a horizontal from it to the right; a long vertical down from the horizontal; then two short horizontals to its right, the upper one first.

\textbf{\zh{心}, four strokes}, at the foot: the left dot, the curve along the foot that hooks up, the middle dot, the right dot.

5 + 4 = nine.

\subsection*{Your turn}
\begin{itemize}
\raggedright
  \item \emph{Write it:} \zh{怎}, all nine, with the model in front of you
  \item \emph{Cover:} write it again from memory, and count the pieces, not the strokes
\end{itemize}

\begin{tcolorbox}[breakable,colback=teal!4,colframe=teal!35!black,title={Before you move on}]
How many strokes in \zh{怎}? (\textbf{Nine}.) Where have you written its top? (\textbf{On the right of \zh{作}}.) What is at the foot? (\textbf{\zh{心}}.)

Source: \href{https://raw.githubusercontent.com/chanind/hanzi-writer-data/68d10a4b21150cae5e1ebbd223eed289cf32d90c/data/\%E6\%80\%8E.json}{Hanzi Writer Data \zh{怎}, pinned stroke paths and medians}.
\end{tcolorbox}
\section[\texorpdfstring{\zh{去}…\zh{怎么走} (qù … zěnme zǒu)}{去…怎么走 (qù … zěnme zǒu)}]{\zh{去}…\zh{怎么走} — how do I get there?}

\label{lesson:ZH-C283-zenme-zou}



\begin{quote}
The word you wrote last, and a verb you have said for a long time.

\begin{itemize}
\raggedright
  \item \emph{Read it:} \zh{怎么}, aloud — \emph{zěnme}
  \item \emph{Recall:} say \emph{zǒu}, to walk — \textbf{R3}, durable
  \item \emph{Recall:} say \zh{火车站} — \emph{huǒchēzhàn}, the station — \textbf{R3}
\end{itemize}
\end{quote}

\begin{grammarlens}[title={Grammar Lens: \zh{去} … \zh{怎么走}, asking the way}]
\begin{quote}
\textbf{\zh{去}, the place, then \zh{怎么走} asks the way there.}
\end{quote}

\begin{quote}
\textbf{\zh{去火车站怎么走}} \emph{qù huǒchēzhàn zěnme zǒu} — How do I get to the station? \textbf{\zh{请问} \zh{去书店怎么走}} \emph{qǐngwèn, qù shūdiàn zěnme zǒu} — Excuse me, how do I get to the bookshop?
\end{quote}

\zh{走} is the \emph{zǒu} you know, to walk; the next lesson puts it in your hand. \zh{请问} in front makes the question polite.

The answer shows the way: \emph{yìzhí zǒu}, straight on; \emph{dào lùkǒu zuǒ guǎi}, turn left at the crossing. zuǒ is the zuǒ of \emph{zuǒbian}; yòu, the yòu of \emph{yòubian}, turns you right.
\end{grammarlens}

\subsection*{Your turn}
Say these aloud:

\begin{itemize}
\raggedright
  \item \zh{去火车站怎么走} — \emph{qù huǒchēzhàn zěnme zǒu} — How do I get to the station?
  \item \zh{请问} \zh{去书店怎么走} — \emph{qǐngwèn, qù shūdiàn zěnme zǒu} — Excuse me, how do I get to the bookshop?
  \item \emph{yìzhí zǒu, dào lùkǒu yòu guǎi} — Straight on, and turn right at the crossing.
  \item ask the way to a place in your own town
\end{itemize}

\begin{tcolorbox}[breakable,colback=teal!4,colframe=teal!35!black,title={Before you move on}]
How do you ask the way to the station? (\textbf{\zh{去火车站怎么走}}.) What does \emph{yìzhí zǒu} tell you? (\textbf{Straight on}.)
\end{tcolorbox}
\section[\texorpdfstring{\zh{走} (zǒu)}{走 (zǒu)}]{\zh{走} — \zh{土}, then four strokes along the foot}

\label{lesson:ZH-W283-zou}



\begin{quote}
The question first, then the character it was waiting for.

\begin{itemize}
\raggedright
  \item \emph{Recall:} say \zh{去火车站怎么走} — \emph{qù huǒchēzhàn zěnme zǒu} — \textbf{R1}, one lesson back, before it decays
\end{itemize}

\begin{quote}
\zh{在} \zh{正}
\end{quote}

\emph{Read it:} each one aloud, with nothing said first — \emph{zài}, \emph{zhèng}
\end{quote}

\subsection*{Script — \zh{土}, then four strokes along the foot}
\hlblockfigure{\detokenize{figures/ZH-W283-zou-filmstrip.pdf}}{How \zh{走} is written, stroke by stroke}

\begin{quote}
\zh{走}
\end{quote}

\textbf{Two pieces, and you write the first already.}

\textbf{\zh{土}, three strokes}, on top, as at the lower right of \zh{在}: a short horizontal, a vertical down through it, and a longer horizontal.

\textbf{The lower piece, four strokes}: a short vertical; a short horizontal to its right, the way the lower part of \zh{正} begins; a left-falling stroke to the lower left; and last the long sweep along the foot to the right.

3 + 4 = seven.

\subsection*{Your turn}
\begin{itemize}
\raggedright
  \item \emph{Write it:} \zh{走}, all seven, with the model in front of you
  \item \emph{Cover:} write it again from memory, and count the pieces, not the strokes
\end{itemize}

\begin{tcolorbox}[breakable,colback=teal!4,colframe=teal!35!black,title={Before you move on}]
How many strokes in \zh{走}? (\textbf{Seven}.) Where have you written its top? (\textbf{At the lower right of \zh{在}}.) Which stroke comes last? (\textbf{The long sweep along the foot}.)

Source: \href{https://raw.githubusercontent.com/chanind/hanzi-writer-data/68d10a4b21150cae5e1ebbd223eed289cf32d90c/data/\%E8\%B5\%B0.json}{Hanzi Writer Data \zh{走}, pinned stroke paths and medians}.
\end{tcolorbox}
\section[wǎng]{wǎng — towards}

\label{lesson:ZH-C283-wang}



\begin{quote}
Before the new one: the question, and the two sides of the street.

\begin{itemize}
\raggedright
  \item \emph{Read it:} \zh{去火车站怎么走}, aloud — \emph{qù huǒchēzhàn zěnme zǒu}
  \item \emph{Recall:} say \emph{zuǒbian}, the left — \textbf{R3}, durable
  \item \emph{Recall:} say \emph{yòubian}, the right — \textbf{R3}
\end{itemize}
\end{quote}

\subsection*{You'll want to know: wǎng}
\emph{wǎng} — \textquotedblleft{}towards; in the direction of\textquotedblright{}.

\begin{quote}
\emph{wǎng qián zǒu} — Go straight ahead.
\end{quote}

\begin{quote}
\emph{dào lùkǒu wǎng zuǒ guǎi} — At the crossing, turn left.
\end{quote}

Its character comes later; for now it is a word for the ear.

\begin{cousinweb}
\emph{wǎng} is inherited from Middle Chinese \emph{hjwangX} (Guangyun, in Baxter's transcription), glossed 'to go' and 'to go to'; its use as 'towards' is not given there.
\end{cousinweb}

\begin{grammarlens}[title={What you've built}]
\emph{wǎng} comes before a direction, and the verb follows: \emph{wǎng qián zǒu}, \emph{wǎng zuǒ guǎi}, \emph{wǎng yòu guǎi}.

qián is the \zh{前} of \zh{前边}: ahead.
\end{grammarlens}

\subsection*{Your turn}
\begin{itemize}
\raggedright
  \item \emph{Say it:} \emph{wǎng}
  \item \emph{Say it:} \emph{wǎng}, once more
  \item \emph{Say it:} \emph{wǎng qián zǒu} — Go straight ahead.
  \item \emph{Say it:} \emph{dào lùkǒu wǎng zuǒ guǎi} — At the crossing, turn left.
  \item \emph{Recall:} say \emph{wǎng yòu guǎi}: turn right
\end{itemize}

\begin{tcolorbox}[breakable,colback=teal!4,colframe=teal!35!black,title={Before you move on}]
What does \emph{wǎng} mean? (\textquotedblleft{}Towards\textquotedblright{}.) What comes after it? (\textbf{A direction, then the verb}.)
\end{tcolorbox}
\section[\texorpdfstring{\zh{再} (zài)}{再 (zài)}]{\zh{再} (zài) — again}

\label{lesson:ZH-C283-zai-again}



\begin{quote}
Before the new one: last lesson's word, and a goodbye you have said from the first chapters.

\begin{itemize}
\raggedright
  \item \emph{Recall:} say \emph{wǎng qián zǒu}, straight ahead — \textbf{R1}, one lesson back
  \item \emph{Recall:} say \zh{再见} — \emph{zàijiàn}, goodbye — \textbf{R3}, durable
\end{itemize}
\end{quote}

\subsection*{You'll want to know: \zh{再}}
\textbf{\zh{再}} — \emph{zài} — \textquotedblleft{}again; once more\textquotedblright{}, for what is still to come.

\begin{quote}
\textbf{\zh{我明天再来}} \emph{wǒ míngtiān zài lái} — I'll come again tomorrow.
\end{quote}

\begin{quote}
\emph{qǐng zài shuō yí cì} — Please say it once more.
\end{quote}

It is the \zh{再} of \zh{再见}: see you again.

\begin{cousinweb}
\textbf{\zh{再}} (\emph{zài}) is inherited from Middle Chinese \emph{tsojH} (Guangyun, in Baxter's transcription), glossed 'twice; a second time'.
\end{cousinweb}

\begin{grammarlens}[title={What you've built}]
\zh{再} goes before the verb, after the person, and it looks forward: the again has not happened yet.

It is not the \zh{在} of \zh{我在家}: the same sound, a different word, a different character.
\end{grammarlens}

\subsection*{Your turn}
\begin{itemize}
\raggedright
  \item \emph{Say it:} \emph{zài}
  \item \emph{Say it:} \emph{zài}, once more
  \item \emph{Say it:} \zh{我明天再来} — \emph{wǒ míngtiān zài lái} — I'll come again tomorrow.
  \item \emph{Say it:} \emph{qǐng zài shuō yí cì} — Please say it once more.
  \item \emph{Recall:} say \zh{再见}, then \zh{我明天再来}
\end{itemize}

\begin{tcolorbox}[breakable,colback=teal!4,colframe=teal!35!black,title={Before you move on}]
What does \zh{再} mean? (\textquotedblleft{}Again\textquotedblright{}, still to come.) Which \zh{在} is it not? (\textbf{The \zh{在} of place, \zh{我在家}}.)
\end{tcolorbox}
\section[\texorpdfstring{\zh{路过} (lùguò)}{路过 (lùguò)}]{\zh{路过} (lùguò) — to pass by}

\label{lesson:ZH-C283-luguo}



\begin{quote}
Before the new one: last lesson's word, then two characters you write.

\begin{itemize}
\raggedright
  \item \emph{Read it:} \zh{我明天再来}, aloud — \emph{wǒ míngtiān zài lái}
\end{itemize}

\begin{quote}
\zh{路} \zh{过}
\end{quote}

\emph{Read it:} each one aloud, with nothing said first — \emph{lù}, \emph{guò}
\end{quote}

\subsection*{You'll want to know: \zh{路过}}
\textbf{\zh{路过}} — \emph{lùguò} — \textquotedblleft{}to pass by\textquotedblright{}, on the way somewhere else.

\begin{quote}
\textbf{\zh{我每天路过学校}} \emph{wǒ měitiān lùguò xuéxiào} — I go past the school every day.
\end{quote}

\begin{quote}
\textbf{\zh{我路过你家}} \emph{wǒ lùguò nǐ jiā} — I'm going past your house.
\end{quote}

Both characters are ones you write.

\begin{cousinweb}
\textbf{\zh{路过}} (\emph{lùguò}) is built from \zh{路} 'road' and \zh{过} 'to pass', both inherited from Middle Chinese, \emph{luH} and \emph{kwaH} (Guangyun, in Baxter's transcription), which glosses the first 'a road' and the second, among other senses, 'to go beyond'.
\end{cousinweb}

\begin{grammarlens}[title={What you've built}]
\zh{路过} is passing a place on your way somewhere else: you do not stop there. \zh{去} is going there.

It is a verb, and the place follows it: \zh{路过学校}.
\end{grammarlens}

\subsection*{Your turn}
\begin{itemize}
\raggedright
  \item \emph{Say it:} \emph{lùguò}
  \item \emph{Say it:} \emph{lùguò}, once more
  \item \emph{Say it:} \zh{我每天路过学校} — \emph{wǒ měitiān lùguò xuéxiào} — I go past the school every day.
  \item \emph{Say it:} \zh{我路过你家} — \emph{wǒ lùguò nǐ jiā} — I'm going past your house.
  \item \emph{Recall:} say \zh{路}, then \zh{过}, then \zh{路过}
\end{itemize}

\begin{tcolorbox}[breakable,colback=teal!4,colframe=teal!35!black,title={Before you move on}]
What does \zh{路过} mean? (\textquotedblleft{}To pass by\textquotedblright{}.) Do you stop at the place? (\textbf{No: it is on the way}.)
\end{tcolorbox}
\section[Practice]{\zh{去书店怎么走}}

\label{lesson:ZH-C283-practice}



\begin{quote}
The chapter closes here, with a stranger asking the way.

\begin{itemize}
\raggedright
  \item \emph{Recall:} say \zh{路过} — \emph{lùguò} — \textbf{R1}, one lesson back
  \item \emph{Recall:} say \zh{再} — \emph{zài}, again — \textbf{R2}
  \item \emph{Recall:} say \emph{wǎng}, towards — \textbf{R2}
\end{itemize}
\end{quote}

\subsection*{The exchange}
\begin{quote}
\textbf{A:} \zh{请问} \zh{去书店怎么走}
\end{quote}

\begin{quote}
\emph{qǐngwèn, qù shūdiàn zěnme zǒu} — Excuse me, how do I get to the bookshop?
\end{quote}

\begin{quote}
\textbf{B:} \emph{wǎng qián zǒu} — Go straight ahead.
\end{quote}

\begin{quote}
\textbf{B:} \emph{dào lùkǒu wǎng yòu guǎi} — Turn right at the crossing.
\end{quote}

\begin{quote}
\textbf{A:} \emph{duìbuqǐ, qǐng zài shuō yí cì} — Sorry, could you say that again?
\end{quote}

\begin{quote}
\textbf{B:} \emph{wǎng qián zǒu, dào lùkǒu wǎng yòu guǎi} — Straight ahead, then right at the crossing.
\end{quote}

\begin{quote}
\textbf{A:} \emph{yuǎn ma} — Is it far?
\end{quote}

\begin{quote}
\textbf{B:} \emph{bù yuǎn, wǒ měitiān lùguò} — Not far. I go past it every day.
\end{quote}

\begin{quote}
\textbf{A:} \zh{谢谢} \zh{再见}
\end{quote}

\begin{quote}
\emph{xièxie, zàijiàn} — Thank you. Goodbye.
\end{quote}

Directions run in the order you walk them, and the person is left out: \emph{wǎng qián zǒu}, then \emph{dào lùkǒu wǎng yòu guǎi}. Asking again is \emph{qǐng zài shuō yí cì}.

\begin{tcolorbox}[breakable,skin=enhanced,colback=orange!6,colframe=orange!45!black,boxrule=0.5pt,arc=1mm,left=6pt,right=6pt,top=4pt,bottom=4pt,fonttitle=\bfseries,title={Read this}]
\emph{Read it:} the exchange through once without stopping, then answer the question below
\end{tcolorbox}

\begin{grammarlens}[title={What you've built — asking the way, and following it}]
A question, three answers, and a way to hear them twice:

\begin{itemize}
\raggedright
  \item \textbf{\zh{去火车站怎么走}} \emph{qù huǒchēzhàn zěnme zǒu} — the question.
  \item \emph{wǎng qián zǒu} — straight ahead.
  \item \emph{dào lùkǒu wǎng zuǒ guǎi} — left at the crossing.
  \item \textbf{\zh{我每天路过}} \emph{wǒ měitiān lùguò} — I pass it every day.
  \item \emph{qǐng zài shuō yí cì} — once more, please.
\end{itemize}

\textbf{Listen.} \emph{Qù xuéxiào zěnme zǒu?}

\textbf{Speak.} Give the way from your door to the nearest shop, in three steps.

\emph{Write it:} \zh{怎么走}
\end{grammarlens}

\subsection*{Your turn}
\begin{itemize}
\raggedright
  \item \emph{Say it:} \zh{去书店怎么走} — \emph{qù shūdiàn zěnme zǒu} — How do I get to the bookshop?
  \item \emph{Say it:} \emph{wǎng qián zǒu, dào lùkǒu wǎng yòu guǎi} — Straight ahead, then right at the crossing.
  \item \emph{Say it:} the whole exchange, both parts
  \item \emph{Write it:} \zh{去火车站怎么走}
\end{itemize}

\begin{tcolorbox}[breakable,colback=teal!4,colframe=teal!35!black,title={Before you move on}]
Ask the way to the bookshop. (\textbf{\zh{去书店怎么走}}.) Say go straight ahead. (\emph{wǎng qián zǒu}.) Ask someone to say it again. (\emph{qǐng zài shuō yí cì}.)
\end{tcolorbox}
